\begin{tikzpicture}[
  x=1cm,
  y=1cm,
  font=\small,
  module/.style={
    draw,
    rounded corners=2pt,
    fill=white,
    align=center,
    minimum height=0.9cm,
    text width=3.35cm,
    inner sep=4pt
  },
  flow/.style={-{Stealth[length=2.3mm]}, thick, draw=black!65},
  trigger/.style={-{Stealth[length=2.3mm]}, dashed, thick, draw=black!60}
]
  \path[draw, rounded corners=4pt, fill=blue!7] (-7.1, 4.7) rectangle (7.1, 6.7);
  \node[anchor=west, font=\bfseries] at (-6.7, 6.15) {表现层};
  \node[module] (ui1) at (-4.3, 5.55) {题目集与题目详情};
  \node[module] (ui2) at (0, 5.55) {代码编辑器与\\统一工作区};
  \node[module] (ui3) at (4.3, 5.55) {提交记录、结果详情\\与 AI 反馈面板};

  \path[draw, rounded corners=4pt, fill=teal!7] (-7.1, 2.0) rectangle (7.1, 4.15);
  \node[anchor=west, font=\bfseries] at (-6.7, 3.6) {业务服务层};
  \node[module, text width=3.2cm] (svc1) at (-4.3, 2.95) {Next.js 页面路由\\与 Server Actions};
  \node[module, text width=3.2cm] (svc2) at (0, 2.95) {认证鉴权、题目读取\\与模板加载服务};
  \node[module, text width=3.2cm] (svc3) at (4.3, 2.95) {提交受理、状态回写\\与结果聚合服务};

  \path[draw, rounded corners=4pt, fill=orange!10] (-7.1, -0.55) rectangle (2.15, 1.35);
  \node[anchor=west, font=\bfseries] at (-6.7, 0.8) {评测执行层};
  \node[module, text width=2.5cm] (judge1) at (-4.7, 0.15) {Submission 创建\\与任务初始化};
  \node[module, text width=2.5cm] (judge2) at (-1.8, 0.15) {编译命令生成\\与日志截断};
  \node[module, text width=2.5cm] (judge3) at (1.0, 0.15) {Docker 容器执行\\与用例判定};

  \path[draw, rounded corners=4pt, fill=green!10] (2.95, -0.55) rectangle (7.1, 1.35);
  \node[anchor=west, font=\bfseries] at (3.25, 0.95) {AI辅助层};
  \node[module, text width=2.75cm] (ai1) at (5.45, -0.05) {异步分析任务、\\DeepSeek 调用\\与结构化反馈生成};

  \path[draw, rounded corners=4pt, fill=gray!10] (-7.1, -3.15) rectangle (7.1, -1.2);
  \node[anchor=west, font=\bfseries] at (-6.7, -1.7) {数据层};
  \node[module, text width=4.0cm] (data1) at (-2.5, -2.2) {Prisma 数据访问层\\User / Problem / Submission / Testcase / CodeAnalysis};
  \node[module, text width=4.0cm] (data2) at (3.2, -2.2) {PostgreSQL 持久化存储\\及 Docker 本地 / 远程运行环境配置};

  \draw[flow] (ui1.south) -- ++(0, -0.35) -| (svc1.north);
  \draw[flow] (ui2.south) -- (svc2.north);
  \draw[flow] (ui3.south) -- ++(0, -0.35) -| (svc3.north);

  \draw[flow] (svc1.south) -- ++(0, -0.3) -| (judge1.north);
  \draw[flow] (svc2.south) -- ++(0, -0.3) -| (judge2.north);
  \draw[flow] (svc3.south) -- ++(0, -0.3) -| (judge3.north);

  \draw[flow] (judge1.south) -- ++(0, -0.35) -| (data1.north);
  \draw[flow] (judge2.south) -- ++(0, -0.35) -| (data1.north);
  \draw[flow] (judge3.south) -- ++(0, -0.35) -| (data2.north);

  \draw[flow] (ai1.south) -- ++(0, -0.35) -| (data2.north);
  \draw[trigger] (judge3.east) -- node[above, font=\scriptsize, fill=white, inner sep=1pt] {判题完成后触发} (ai1.west);
  \draw[flow] (svc3.east) -- ++(0.55, 0) |- (ai1.north);
\end{tikzpicture}
